\section{User interfaces}\label{sec:interfaces}

QtRocket presents three user-facing surfaces, all built over the same Qt-free core
(\cref{sec:overview}) but at very different levels of maturity. The Qt6 Widgets application
\code{qtrocket} is a \emph{motor-and-launch} front end: it imports motors, stages a point-mass
launch, flies it, and plots the result --- but it cannot author, open, or save a rocket design.
The headless command-line interface (CLI) \code{qtrocket-cli} is the most complete surface and
the only one that can author designs: a line-oriented read-eval-print loop (REPL) whose
stdout doubles as a machine-parseable protocol, making it scriptable from a pipe or a test
harness. The standalone \code{qtrocket-visualizer} is a self-contained OpenGL viewer for saved
\code{.qrd} design files that positions every rendered part with the \emph{same} placement
resolver the simulator's mass pass consumes (\cref{sec:placement}). This section walks each
surface in turn and closes with the asymmetry between them;
\cref{fig:interfaces-components} maps the classes involved and the core services each
surface actually drives.

\tikzfig{component-interfaces}{The three front ends and the core services each one drives.
Solid edges are the load-bearing call paths at the pinned commit; the GUI's missing edge to
the design serializer --- its File actions ship disabled --- is the maturity asymmetry this
section documents.}{fig:interfaces-components}

\subsection{The Qt6 Widgets GUI}\label{sec:interfaces-gui}

The \srcf{gui/} directory is the only Qt-dependent code in the main tree --- roughly a
thousand lines of hand-written widget code compiled directly into the \code{qtrocket}
executable (\src{CMakeLists.txt}{159}), plus the vendored 43{,}000-line QCustomPlot 2.1.1
plotting library (\src{gui/qcustomplot.h}{23}), which is compiled with \code{-w} so the rest
of the project can keep \code{-Werror} (\src{CMakeLists.txt}{152}) and is excluded from
coverage metrics (\src{CMakeLists.txt}{262}). No GUI source is ever built into a library or a
test binary; no test links \code{Qt6::Widgets} (\cref{sec:testing}).

\subsubsection{The entry point: \code{QApplication} on a spawned thread}
\label{sec:interfaces-gui-thread}

\code{main()} creates the \code{Logger} and the \code{QtRocket} controller singleton on the
process main thread --- no Qt yet (\src{main.cpp}{20}) --- and then blocks in
\cppinline|gui::run(qtrocket, argc, argv)| (\src{main.cpp}{25}). The entire GUI entry point is
seven lines (\cref{lst:interfaces-guirun}).

\begin{listing}[htbp]
\begin{minted}{cpp}
int run(QtRocket* qtRocket, int argc, char* argv[])
{
   int ret = 0;
   std::thread guiThread(guiWorker, qtRocket, argc, argv, std::ref(ret));
   guiThread.join();
   return ret;
}
\end{minted}
\caption{The whole GUI entry point: spawn a worker thread, join it immediately, and return
the exit code through an out-parameter. \code{QApplication} is constructed inside
\code{guiWorker}, on the spawned thread. \srccap{gui/GuiRunner.cpp}{59}}
\label{lst:interfaces-guirun}
\end{listing}

\code{guiWorker} constructs \code{QApplication} on that spawned thread
(\src{gui/GuiRunner.cpp}{30}), shows the \code{MainWindow}, and runs the entire Qt event loop
there (\src{gui/GuiRunner.cpp}{51}); the exit code travels back through an \cppinline|int&|
passed via \cppinline|std::ref| (\src{gui/GuiRunner.cpp}{62}). The reason is layering, not
concurrency, and the header says so: the function ``lives in the gui layer so the QtRocket
controller stays Qt-free and links into the headless CLI'' (\src{gui/GuiRunner.h}{11}).
Keeping every Qt construction out of \code{main()}'s call chain is what lets the same
controller object serve \code{qtrocket-cli} without a display.

\begin{keyinsight}[Qt's GUI thread is not the process main thread]
Qt treats whichever thread constructs \code{QCoreApplication} as its main/GUI thread; the
documented convention is to construct it in \code{main()}'s thread. QtRocket inverts this:
the spawned thread owns the event loop (\src{gui/GuiRunner.cpp}{30}) while the true main
thread does nothing but sleep in \code{join()} (\src{gui/GuiRunner.cpp}{63}). Because the
join is immediate, there is no concurrency benefit --- the process is behaviorally
single-threaded-at-a-time. The pattern buys layering purity and imports a portability risk:
it works on Linux/xcb, but platforms whose native event loop must own the first thread
(macOS/Cocoa) do not support a GUI on a secondary thread.
\end{keyinsight}

Two smaller facts about \code{guiWorker} are worth recording. The application-level icon is
set from \code{:/qtrocket.png} (\src{gui/GuiRunner.cpp}{31}), a path absent from the
resource file --- the only embedded image is \code{:/images/resources/rocket64.png}
(\src{qtrocket.qrc}{4}), which the window-level icon in the \code{.ui} correctly uses
(\src{gui/MainWindow.ui}{16}) --- so the application icon is silently null. And the
translation loop tries to load \code{:/i18n/} resources the \code{qrc} does not contain
(\src{gui/GuiRunner.cpp}{40}, \src{qtrocket.qrc}{1}): internationalization is dead plumbing
at this commit (\cref{sec:incomplete}).

\subsubsection{MainWindow: the shell, the disabled File menu, and the stub panes}
\label{sec:interfaces-gui-mainwindow}

\code{MainWindow}'s constructor builds the \code{.ui} skeleton and adds its two working tabs
programmatically: \code{CannonballTab} and \code{SimOptionsTab}
(\src{gui/MainWindow.cpp}{30}). The skeleton around them is a vertical splitter holding a
horizontal splitter --- a \code{RocketTreeView} beside the tab widget
(\src{gui/MainWindow.ui}{49}) --- above a \code{RocketModelerView}
(\src{gui/MainWindow.ui}{66}). Both custom widgets are Qt Designer promoted classes
(\src{gui/MainWindow.ui}{187}), and both are verified constructor stubs: each \code{.cpp} is
six lines with an empty constructor body (\src{gui/RocketTreeView.cpp}{3},
\src{gui/RocketModelerView.cpp}{3}). No \code{setModel} call exists anywhere in \srcf{gui/},
so the tree view renders an empty pane and the modeler view a blank widget. More broadly,
nothing in the GUI observes part-tree edits at all: a repository-wide search for any
structure-changed callback --- and for callback members in \srcf{QtRocket.h} or
\srcf{model/RocketModel.h} --- comes back empty at this commit. The rocket-structure UI is
a frame with nothing wired behind it (\cref{sec:incomplete}).

The File menu makes the same point from the other direction. It lists New, Open, Save, Save
As, Close, and Quit (\src{gui/MainWindow.ui}{82}), but five of the six ship permanently
disabled in the \code{.ui} --- \code{actionNew} (\src{gui/MainWindow.ui}{117}),
\code{actionOpen} (\src{gui/MainWindow.ui}{129}), \code{actionSave}
(\src{gui/MainWindow.ui}{141}), \code{actionClose} (\src{gui/MainWindow.ui}{153}), and
\code{actionSave\_As} (\src{gui/MainWindow.ui}{170}) --- and none has a slot. Only Quit is
live, wired to \code{close()} (\src{gui/MainWindow.cpp}{102}). The consequence is stated
plainly: \emph{the GUI cannot open a \code{.qrd} design}. The shared
\code{model::DesignSerializer} (\cref{sec:persistence}) is consumed by the CLI
(\src{cli/Repl.cpp}{957}) and the visualizer (\src{visualizer/VisualizerWindow.cpp}{184}), and
a search of \srcf{gui/} for any design-persistence reference returns nothing.

The two menus that do work are small: Tools\,$\to$\,Save Motor Database opens a
\code{*.qmd} save dialog (start directory hard-coded to the literal \code{/home},
\src{gui/MainWindow.cpp}{76}), appends the \code{.qmd} suffix when missing
(\src{gui/MainWindow.cpp}{84}), and calls \code{saveMotorDatabase} inside a try/catch that
surfaces failures in a message box (\src{gui/MainWindow.cpp}{91}); Help\,$\to$\,About shows
a modal copyright dialog (\src{gui/MainWindow.cpp}{66}).

\subsubsection{CannonballTab and SimOptionsTab}\label{sec:interfaces-gui-tabs}

\textbf{CannonballTab} is the launch panel for the point-mass (``cannonball'') flight model:
the rocket is a mass with one drag coefficient and one reference area, so the tab exposes
exactly five numbers and a motor. The line edits and their shipped defaults are: initial
velocity 5.0\,m/s (\src{gui/CannonballTab.ui}{29}), launch angle 0.0$^\circ$
(\src{gui/CannonballTab.ui}{46}), mass 1.0\,kg (\src{gui/CannonballTab.ui}{60}), drag
coefficient 0.5 (\src{gui/CannonballTab.ui}{74}), and reference area
0.001134\,$\mathrm{m^2}$ --- the cross-section of a 38\,mm-diameter airframe
(\src{gui/CannonballTab.ui}{91}). Only two fields carry validators: the angle is clamped to
$[0,90]$ degrees measured \emph{from vertical}, so the default $0^\circ$ means straight up
(\src{gui/CannonballTab.cpp}{34}), and the area's validator bottom is zero
(\src{gui/CannonballTab.cpp}{39}). The other three parse through
\code{QString::toDouble()}, which silently yields 0.0 on junk; the model-side mass setter then
silently ignores non-positive values (\src{model/RocketModel.h}{87}) and the area setter
negative ones (\src{model/RocketModel.h}{74}), so the display and the model can disagree with
no feedback.

Calculate Trajectory reads the five fields, decomposes the initial speed by the
angle-from-vertical convention,
\begin{equation}\label{eq:interfaces-angle}
v_x = v_0\sin\theta, \qquad v_y = 0, \qquad v_z = v_0\cos\theta,
\end{equation}
with degrees converted through the rounded constant 57.2958 $\approx 180/\pi$
(\src{gui/CannonballTab.cpp}{105}), builds the initial \code{StateData} at the origin
(\src{gui/CannonballTab.cpp}{107}), pushes mass, drag coefficient, and reference area into the
\code{RocketModel} (\src{gui/CannonballTab.cpp}{111}), and then calls
\code{setInitialState} followed by \code{launchRocket}
(\src{gui/CannonballTab.cpp}{115}). The launch is fully synchronous: the entire ODE
integration (\cref{sec:simcore}) runs inside the button slot
(\src{QtRocket.cpp}{54}), so the window is frozen for the duration of the flight. On return
the tab opens the \code{AnalysisWindow}; the \code{setModal(false)} preceding
\code{exec()} is ineffective, since \code{QDialog::exec()} always runs application-modal
(\src{gui/CannonballTab.cpp}{118}).

The tab's enable logic is deliberately centralized: the launch button is enabled iff the
rocket has a motor, \cppinline|isMotorSet()| (\src{gui/CannonballTab.cpp}{85}, contract at
\src{gui/CannonballTab.h}{47}), and every motor-selection path re-invokes the same helper
rather than toggling the button ad hoc (\src{gui/CannonballTab.cpp}{77},
\src{gui/CannonballTab.cpp}{170}, \src{gui/CannonballTab.cpp}{241}). The motor paths all
route through the \code{MotorModelDatabase} (\cref{sec:motors}): a file loader that
dispatches on extension (\code{.eng} to the RASP importer, everything else to the RockSim
\code{.rse} importer, \src{gui/CannonballTab.cpp}{136}), \code{.qmd} database load/save
dialogs (\src{gui/CannonballTab.cpp}{173}), and a Set Motor button that resolves the combo's
common name against the database (\src{gui/CannonballTab.cpp}{231}). The engine combo is
always rebuilt from \code{listMotors()} and cleared first --- the database is the single
source of truth for its contents (\src{gui/CannonballTab.cpp}{153}, contract at
\src{gui/CannonballTab.h}{44}).

\textbf{SimOptionsTab} applies every field live --- there is no Ok/Apply button
(\src{gui/SimOptionsTab.h}{19}). The timestep validator is strictly positive because a
non-positive $dt$ would hang the run loop (\src{gui/SimOptionsTab.cpp}{25}). The three
combos are populated from the model registries --- the source of truth for valid names:
atmosphere \{Constant Atmosphere, US Standard 1976, Vacuum\} and gravity \{Constant
Gravity, Spherical Gravity\} from \code{sim::Environment} (\src{sim/Environment.h}{103};
\cref{sec:environment}), integrator \{Runge-Kutta 4th Order, Runge-Kutta-Fehlberg\} from
\code{sim::Integrator} (\src{sim/Integrator.h}{32}; \cref{sec:simcore}). Signals are
connected only \emph{after} population so the \code{addItem} calls do not fire spurious
applies, and one apply of all four displayed values is forced at construction so the engine
always matches the display (\src{gui/SimOptionsTab.cpp}{48},
\src{gui/SimOptionsTab.cpp}{69}). The atmosphere and gravity slots mutate the shared
\code{Environment} \emph{in place} --- changing the atmosphere model must not clobber
gravity, and the rest of the application holds the same \code{shared\_ptr}
(\src{gui/SimOptionsTab.cpp}{89}); timestep and integrator forward through the controller to
the propagator (\src{gui/SimOptionsTab.cpp}{101}, \src{QtRocket.h}{36}).

\subsubsection{Motor search over the network}\label{sec:interfaces-gui-motorsearch}

\code{ThrustCurveMotorSelector} is the dialog behind the Get Thrustcurve Motor Data button:
three facet combos (manufacturer, diameter, impulse class), a search button, a motor list, and
an embedded QCustomPlot thrust-curve preview
(\src{gui/ThrustCurveMotorSelector.ui}{71}). Get Metadata fetches the search facets from
thrustcurve.org and appends them to the combos (\src{gui/ThrustCurveMotorSelector.cpp}{48});
the combos are never cleared first, so repeated clicks duplicate every entry, and the
manufacturer facet displays the API's short \emph{code} rather than the display name because
the search API keys on the code (\src{model/MotorModelDatabase.cpp}{151}). Search builds a
\code{MotorQuery} from the non-empty combos only, leaving unset facets unconstrained
(\src{gui/ThrustCurveMotorSelector.cpp}{63}); every result is merged into the local database
(\src{model/MotorModelDatabase.cpp}{171}), which is why Set Motor can later resolve it by
common name and plot its curve (\src{gui/ThrustCurveMotorSelector.cpp}{83}).

The network path beneath those slots is fully synchronous, end to end, on the GUI thread.
The dialog calls \code{MotorModelDatabase}, which lazily constructs the real
\code{ThrustCurveClient} behind the \code{ThrustCurveAPI} interface
(\src{model/MotorModelDatabase.cpp}{137}; \cref{sec:motors}); the client issues HTTP GETs
through \code{utils::CurlConnection::get} (\src{model/ThrustCurveClient.cpp}{270},
\src{model/ThrustCurveClient.cpp}{301}), which is a blocking
\code{curl\_easy\_perform} with a 10\,s connect timeout and a 30\,s total timeout
(\src{utils/CurlConnection.cpp}{54}, \src{utils/CurlConnection.cpp}{58}).

\begin{keyinsight}[One click, up to 101 blocking round trips]
\code{searchMotors} fans out: after the search response it loops the results and fetches each
motor's thrust-curve samples with a separate \code{download.json} request
(\src{model/ThrustCurveClient.cpp}{320}), and the default result cap is
\code{maxResults=100} (\src{model/ThrustCurveClient.cpp}{298}). One Search click is therefore
$1 + N$ blocking HTTP round trips with $N \le 100$, each bounded only by the 30\,s libcurl
timeout, all executed in the slot body on the Qt event-loop thread. The code states this
plainly on both sides of the seam: ``each result costs one download.json round trip below,
performed synchronously on the caller's (often GUI) thread''
(\src{model/ThrustCurveClient.cpp}{293}), and ``these calls run synchronously on the caller's
thread (including the GUI thread), so a hung connection must time out rather than wedge the
app'' (\src{utils/CurlConnection.cpp}{52}). There is no \code{QThread},
\code{QtConcurrent}, or \code{QNetworkAccessManager} anywhere in \srcf{gui/}: the timeouts
are the only thing standing between a slow server and a multi-minute frozen window.
\end{keyinsight}

\subsubsection{AnalysisWindow: exactly three fixed plots}\label{sec:interfaces-gui-analysis}

\code{AnalysisWindow} is the post-flight view: three buttons drawing into a single promoted
QCustomPlot widget (\src{gui/AnalysisWindow.ui}{16}), each clearing the previous graph and
enabling drag/zoom. Plot Altitude reads the recorded $(t, \code{StateData})$ history from
\code{QtRocket::getStates()} (\src{gui/AnalysisWindow.cpp}{39}) and plots
\cppinline|position[2]| --- altitude $z$ in meters --- against time
(\src{gui/AnalysisWindow.cpp}{49}). Plot Velocity plots \cppinline|velocity[2]|, the vertical
velocity component, not speed (\src{gui/AnalysisWindow.cpp}{74}). Motor Curve plots the
motor's thrust samples in Newtons against time (\src{gui/AnalysisWindow.cpp}{86}). These
three are the entire analysis surface: the plotted variables are hard-coded, and there is no
downrange, speed, acceleration, or mass plot. Axis ranges are spanned from the data
(\src{gui/AnalysisWindow.cpp}{55}); the labels are minimal (the altitude axis is labeled
literally ``Z'', \src{gui/AnalysisWindow.cpp}{54}). Since only the three linear degrees of
freedom are integrated at this commit (\cref{sec:simcore}), the plot set matches what the
state actually contains --- but the CSV the CLI exports carries strictly more
(\cref{sec:interfaces-cli-launch}).

\subsection{The command-line REPL}\label{sec:interfaces-cli}

\code{qtrocket-cli} is the headless front end: a 26-line \code{main}
(\srcf{cli/CliMain.cpp}) wrapping \code{cli::Repl} (\srcf{cli/Repl.cpp}, 985 lines), a
line-oriented REPL over the same \code{QtRocket} API the GUI drives. It reads line by line,
so the identical binary works interactively, from a pipe, or from a redirected script
(\src{cli/Repl.h}{22}). The Repl is compiled into its own \code{cli} static library
(\src{CMakeLists.txt}{235}) so that the test suites drive the exact object code a user types
into (\cref{sec:testing}); the executable links that library and nothing else
(\src{CMakeLists.txt}{297}).

\subsubsection{Process model and the stdout protocol}\label{sec:interfaces-cli-protocol}

\begin{listing}[htbp]
\begin{minted}{cpp}
int main(int /*argc*/, char* /*argv*/[])
{
   // Logger shares stdout with the CLI's machine-readable output, so run at ERROR. The "[ERROR]"
   // prefix keeps those lines distinct from the CLI's own "OK"/"ERR"/"CSV".
   utils::Logger::getInstance()->setLogLevel(utils::Logger::ERROR_);
\end{minted}
\caption{The CLI's first act is pinning the Logger to ERROR: stdout is a protocol surface,
and log chatter would corrupt it. The rest of the 26-line \code{main} prints the
\code{\#}-prefixed banner and hands \code{std::cin}/\code{std::cout} to the Repl
(\srccap{cli/CliMain.cpp}{13}).}
\label{lst:interfaces-climain}
\end{listing}

The stdout contract is simple and consistently enforced. Every command's first response line
is machine-parseable: \code{OK <cmd>[: detail]} on success or \code{ERR <cmd>: reason} (or
\code{ERR usage: ...}) on failure. Payload rows --- motor listings, model-name lists, CSV
rows, the part tree, the \code{status} block --- follow the \code{OK} line unprefixed.
Human commentary is \code{\#}-prefixed: the startup banner
(\src{cli/CliMain.cpp}{21}) and the entire \code{help} text, which itself ends with
\code{OK help} (\src{cli/Repl.cpp}{259}). Input lines beginning with \code{\#} are treated
as script comments (\src{cli/Repl.cpp}{256}). Two extra line types exist only for
\code{launch}: a \code{WARN:} line on non-nominal termination (\src{cli/Repl.cpp}{737}) and a
\code{CSV <absolute-path>} line naming the flight artifact (\src{cli/Repl.cpp}{740}). This is
why the Logger is pinned to ERROR before anything else runs
(\src{cli/CliMain.cpp}{15}): the protocol and the log share one stream, so the log must be
nearly silent and visually distinct when it does speak.

Parsing is equally deliberate. Only the command word is lowercased --- arguments keep their
case so motor codes like \code{G80T} and file paths survive (\src{cli/Repl.cpp}{252}).
Numeric arguments parse strictly: trailing garbage (e.g.\ \code{0.5abc}) rejects the token
rather than truncating it (\src{cli/Repl.cpp}{119}). An unknown \emph{command} is an
\code{ERR} (\src{cli/Repl.cpp}{981}), but an unknown \emph{design key} in
\code{newdesign}/\code{addpart} warns and continues, so a newer file or CLI's extra keys do
not break an older one (\src{cli/Repl.cpp}{205}).

\begin{keyinsight}[Automation must parse \code{ERR}, not \code{\$?}]
The Repl's run loop reads until EOF or \code{quit}/\code{exit} and then unconditionally
returns 0 (\src{cli/Repl.cpp}{232}), which \code{main} returns as the process exit code
(\src{cli/CliMain.cpp}{25}). A script consisting entirely of failing commands still exits 0.
Pipelines therefore cannot use the shell exit status to detect failure; they must scan stdout
for \code{ERR} lines. One edge case complicates even that: \code{launch} is the one command
that can emit both --- if the hardcoded CSV artifact cannot be opened, an
\code{ERR launch: cannot open ...} line prints and the \code{OK launch:} summary still
follows (\src{cli/Repl.cpp}{717}), so a strict ``first line decides'' parser must treat
\code{launch} specially.
\end{keyinsight}

\subsubsection{The command inventory}\label{sec:interfaces-cli-commands}

The dispatch chain in \code{Repl::execute} (\src{cli/Repl.cpp}{245}) contains 32 handler
branches accepting 33 command words --- \code{quit} and \code{exit} are aliases sharing one
branch (\src{cli/Repl.cpp}{975}). \Cref{tab:interfaces-commands} is the complete inventory,
grouped by function. Two session contracts pinned in the \code{help} text apply across the
groups: part ids reset on reload, and the launch/atmosphere settings are session
configuration, never saved in a design file (\src{cli/Repl.cpp}{294}); and drag uses the
active atmosphere's density, so selecting the Vacuum atmosphere reduces the force model to
thrust plus gravity (\src{cli/Repl.cpp}{298}).

\begingroup\small
\begin{longtable}{@{}>{\ttfamily}p{4.6cm}>{\raggedright\arraybackslash}p{9.0cm}r@{}}
\caption{The complete \code{qtrocket-cli} command inventory at commit \code{254cd41}: 33
command words across 32 handler branches. Angle brackets mark required arguments, square
brackets optional ones. The right column gives the opening line of each command's handler
branch in \srcfcap{cli/Repl.cpp}.}\label{tab:interfaces-commands}\\
\toprule
\textrm{\textbf{Command and arguments}} & \textbf{Semantics} & \textbf{Line}\\
\midrule
\endfirsthead
\multicolumn{3}{@{}l}{\small\emph{Table \ref{tab:interfaces-commands}, continued}}\\
\toprule
\textrm{\textbf{Command and arguments}} & \textbf{Semantics} & \textbf{Line}\\
\midrule
\endhead
\midrule
\multicolumn{3}{r@{}}{\small\emph{continued on next page}}\\
\endfoot
\bottomrule
\endlastfoot
\multicolumn{3}{@{}l}{\textbf{Session and output}}\\
\addlinespace[2pt]
help & Prints the \code{\#}-prefixed command list, then \code{OK help}. & 259\\
status & Dumps the staged configuration: motor, dry mass (kg), drag coefficient, reference
  area ($\mathrm{m^2}$), velocity (m/s), angle (deg from vertical), atmosphere, gravity,
  integrator, database size. & 652\\
states [stride] & Prints the last run's state series as CSV to stdout, keeping every
  stride-th sample (clamped $\ge 1$); \code{ERR} when no run exists. & 744\\
save <path.csv> & Writes the full series (stride 1) to a caller-chosen file; echoes the
  absolute path. & 759\\
quit | exit & Prints \code{OK bye} and ends the session. & 975\\
\addlinespace
\multicolumn{3}{@{}l}{\textbf{Motor database}}\\
\addlinespace[2pt]
loadmotors <file.rse|file.eng> & Imports a motor file; \code{.eng} dispatches to the RASP
  importer, anything else to the RockSim \code{.rse} importer; reports the count added. & 303\\
savedb <file.qmd> & Saves the motor database as XML. & 327\\
loaddb <file.qmd> & Additive database load; reports new and total motor counts. & 348\\
tcfacets & Fetches the thrustcurve.org search facets (manufacturers, diameters, impulse
  classes). & 371\\
tcsearch <k=v>... & Online search; keys \code{manufacturer}, \code{diameter},
  \code{impulseClass}. Results merge into the local database. & 392\\
listmotors [substr] & Lists common names with average thrust (N) and total impulse
  ($\mathrm{N\,s}$); optional substring filter. & 429\\
setmotor <code> & Exact common-name lookup; installs the motor on the rocket and stages the
  session mirror. & 453\\
\addlinespace
\multicolumn{3}{@{}l}{\textbf{Staged flight configuration}}\\
\addlinespace[2pt]
setmass <kg> & Must be $> 0$; applied immediately as a direct mass override on the root
  part. & 479\\
setdrag <cd> & Drag coefficient; no range check; applied immediately. & 497\\
setarea <m\textasciicircum 2> & Must be $\ge 0$; applied immediately and sets the manual-override flag. & 510\\
setvelocity <m/s> & Staged in the Repl; consumed at \code{launch}. & 528\\
setangle <deg> & Degrees from vertical ($0$ = straight up); staged. & 540\\
settimestep <s> & Must be $> 0$; forwarded to the propagator. & 552\\
listatmospheres & Lists the atmosphere-model registry names. & 569\\
setatmosphere <name> & Validates the (possibly multi-word) name against the registry, then
  applies it to the shared environment. & 577\\
listgravity & Lists the gravity-model registry names. & 596\\
setgravity <name> & Validated against the registry; applied in place. & 604\\
listintegrators & Lists the integrator registry names. & 623\\
setintegrator <name> & Validated; forwarded to the propagator. & 632\\
\addlinespace
\multicolumn{3}{@{}l}{\textbf{Simulation control}}\\
\addlinespace[2pt]
launch & Flies the staged configuration; prints a summary block, an optional \code{WARN:}
  line, and a \code{CSV <path>} artifact line (\cref{sec:interfaces-cli-launch}). & 671\\
\addlinespace
\multicolumn{3}{@{}l}{\textbf{Design authoring}}\\
\addlinespace[2pt]
newdesign <type> [k=v...] & Builds a new root part via the factory and installs it with
  \code{setRoot}; drops any staged motor. & 784\\
addpart <parentId|root> <type> [k=v...] [z=<m>] & Builds a child via the factory and
  attaches it to the parent with \code{abut(z)}. & 825\\
listparts & Prints the indented part tree (one \code{[id] type "name" m=<kg>} line per
  part) plus a composite mass/CG footer at $t = 0$. & 880\\
removepart <id> & Removes a subtree by id; refuses the root; re-checks the motor, which may
  have been in the removed subtree. & 894\\
cleardesign & Resets the rocket to the boot placeholder part and clears the staged motor. & 817\\
savedesign <file.qrd> & Writes the design via \code{DesignSerializer}
  (\cref{sec:persistence}). & 926\\
loaddesign <file.qrd> & Atomic load via \code{DesignSerializer}; re-syncs the staged
  mirrors (drag, area, motor) from the loaded rocket. & 946\\
\end{longtable}
\endgroup

\subsubsection{Launching and the CSV artifact}\label{sec:interfaces-cli-launch}

\code{launch} guards on a staged motor (\src{cli/Repl.cpp}{673}), decomposes the staged speed
and angle exactly as the GUI does --- \cref{eq:interfaces-angle}, with the same rounded
57.2958 constant (\src{cli/Repl.cpp}{35}) --- and then calls the same two controller entry
points, \code{setInitialState} and \code{launchRocket} (\src{cli/Repl.cpp}{687}). The summary
statistics are read back from the propagator's live-tracked trajectory statistics rather than
re-derived (\src{cli/Repl.cpp}{701}), and the summary block prints steps, final time, apogee,
maximum speed, downrange distance $\sqrt{x^2+y^2}$, and landing position at fixed
three-decimal precision (\src{cli/Repl.cpp}{728}).

Every successful launch also writes a hardcoded side effect: the full state series goes to
\code{qtrocket\_run.csv} in the current working directory
(\src{cli/Repl.cpp}{712}), and the \code{CSV} protocol line reports its absolute path. One
writer serves the \code{launch} artifact, the \code{states} stdout dump, and the
\code{save} export (\src{cli/Repl.cpp}{102}), so all three emit the identical schema at
nine significant digits: \code{t,x,y,z,vx,vy,vz,mass,cg\_x,cg\_y,cg\_z,Ixx,Iyy,Izz} ---
time in seconds, position in meters, velocity in m/s, mass in kg, center of gravity in
meters, and the inertia-tensor diagonal in $\mathrm{kg\,m^2}$, all SI
(\src{cli/Repl.cpp}{104}). Note that the CSV already carries the time-varying mass, CG, and
inertia diagonal that the GUI's three plots never show.

Non-nominal terminations map to text through an exhaustive switch with no default case, so a
new \code{TerminationReason} enumerator fails to compile in the CLI until it is handled
(\src{cli/Repl.cpp}{40}); the \code{NoLiftoff} message even carries its diagnostic criterion
in prose --- ``never cleared 1 m after 3 s (check thrust vs weight)''
(\src{cli/Repl.cpp}{46}).

\subsubsection{Design authoring from the REPL}\label{sec:interfaces-cli-design}

The design commands make the CLI the only surface that can author a rocket. Geometry
arrives as whitespace-delimited \code{key=value} tokens parsed into the shared
\code{PartParams} struct --- the same vocabulary the factory, the serializer, and the disk
format use, so field names cannot drift between front ends (\cref{sec:parts});
\code{newdesign} and \code{addpart} then drive \code{model::part::makePart}, whose
exceptions (missing required field, unknown type) surface directly as \code{ERR} lines
(\src{cli/Repl.cpp}{803}, \src{cli/Repl.cpp}{863}).

Attachment, however, has exactly one spelling. \code{addpart} always seats the child with
\cppinline|model::part::abut(offset.z())| --- the child abuts its parent's aft plane with an
optional forward standoff \code{z} (\src{cli/Repl.cpp}{872}); the \code{x=}/\code{y=} tokens
are parsed into a full vector and then dropped, since off-axis placement is meaningless in
3-DOF (\src{cli/Repl.cpp}{189}). The other two placement verbs, \code{nestInBore} and
\code{seatOnWall}, exist in the model layer (\src{model/parts/Placement.h}{134}) and
round-trip through the file format, but have no REPL spelling: the \code{.qrd} format is
ahead of the CLI's authoring vocabulary (\cref{sec:persistence}).

\code{savedesign} and \code{loaddesign} are thin wrappers over the shared
\code{DesignSerializer} (\src{cli/Repl.cpp}{936}, \src{cli/Repl.cpp}{957});
\cref{sec:persistence} covers the format and its fail-safe atomic load. After a successful
load the Repl re-syncs its staged mirrors \emph{from} the rocket --- drag coefficient,
reference area, and the motor name (\src{cli/Repl.cpp}{964}); the displayed dry mass is not
re-synced, one of several small session/model seams inventoried in \cref{sec:incomplete}.
Structural edits keep the motor sound throughout: \code{newdesign} clears the staged motor
because \code{setRoot} drops it (\src{cli/Repl.cpp}{812}), and \code{removepart} re-checks
\code{isMotorSet()} because the motor may have been in the removed subtree
(\src{cli/Repl.cpp}{921}).

\subsection{The standalone visualizer}\label{sec:interfaces-viz}

\code{qtrocket-visualizer} is a separate executable under \srcf{visualizer/} (namespace
\code{viz}): a Qt6 OpenGL viewer for \code{.qrd} design files that links
\code{Qt6::Widgets}, \code{Qt6::OpenGLWidgets}, \code{Qt6::OpenGL}, and
\code{qtrocket\_core} --- and, by stated design, never touches \srcf{gui/} or \srcf{cli/}
code (\src{visualizer/CMakeLists.txt}{27}, \src{visualizer/CMakeLists.txt}{5}). It shares
exactly two things with the rest of the system, and they are the two that matter: the design
loader and the placement resolver.

\subsubsection{Startup and the GL-context workaround}\label{sec:interfaces-viz-startup}

The 41-line \code{main} does three deliberate things. It disables
Qt's rendering-hardware-interface (RHI) widget backing store, which on EGL/NVIDIA otherwise
tries and fails to create a GLES2 context before the viewport exists
(\src{visualizer/main.cpp}{22}). It installs a minimal default surface format --- a 24-bit
depth buffer and nothing else: no OpenGL version, no core profile, no multisampling
(\src{visualizer/main.cpp}{24}). And it opens an optional \code{.qrd} path from
\code{argv[1]} (\src{visualizer/main.cpp}{35}). The empty format request is the point, not an
omission: forcing a 3.3 core profile as the application-wide default poisons every context Qt
creates, and on EGL (Wayland plus the NVIDIA driver) the driver then has no matching window
\code{EGLConfig} --- context creation fails with \code{EGL\_BAD\_MATCH} and the 3D view goes
black, while GLX honors the same request, so the app would run on one machine and not another
(\src{visualizer/main.cpp}{8}).

Accepting whatever context the platform grants has two consequences the renderer carries
throughout. Every shader ships as a twin pair --- desktop GLSL~1.20 and OpenGL~ES~1.00
(GLSL is the OpenGL Shading Language), written in the common
\code{attribute}/\code{varying} subset --- selected at initialization by
\cppinline|context()->isOpenGLES()| (\src{visualizer/RocketGLWidget.cpp}{23},
\src{visualizer/RocketGLWidget.cpp}{314}), with attribute locations bound by name before
linking since neither dialect has layout qualifiers
(\src{visualizer/RocketGLWidget.cpp}{201}). And wireframe mode is implemented by expanding
every triangle into its three edges as explicit \code{GL\_LINES} vertices, because
\code{glPolygonMode} does not exist on GLES2 (\src{visualizer/RocketGLWidget.cpp}{518}).

\subsubsection{From \code{.qrd} to pixels: the shared loader and the shared resolver}
\label{sec:interfaces-viz-mesh}

\code{VisualizerWindow} owns a headless \code{model::RocketModel} and a deliberately
\emph{empty} \code{MotorModelDatabase}: a design whose motor cannot be re-resolved loads
anyway, with the miss reduced to a harmless warning
(\src{visualizer/VisualizerWindow.h}{29}). \code{openFile} delegates to the same static
\code{model::DesignSerializer::load} the CLI uses (\src{visualizer/VisualizerWindow.cpp}{184};
\cref{sec:persistence}) inside a try/catch --- a failed load pops a warning box and leaves
any currently displayed model in place (\src{visualizer/VisualizerWindow.h}{41}). On success
the window rebuilds the render list via \code{viz::buildRocketMeshes} and updates the status
bar with the file, the rocket name, and a recursive part count
(\src{visualizer/VisualizerWindow.cpp}{284}). The side panel offers four preset color
schemes, a per-type color picker, and grid/axes/wireframe toggles wired directly to the
viewport's slots (\src{visualizer/VisualizerWindow.cpp}{105}).

\code{RocketMesh} turns the part tree into triangles. Each part type maps to a primitive
builder --- cone, tube (solid or annular), UV sphere, or fin prism set --- via
\code{dynamic\_cast} dispatch; a \code{Motor} or unknown type produces no geometry
(\src{visualizer/RocketMesh.cpp}{450}). All primitives are built in local coordinates
centered on the origin with $+z$ forward (\src{visualizer/RocketMesh.h}{70}), so positioning
is a pure translation. The positions themselves come from the placement resolver
(\cref{lst:interfaces-resolver}).

\begin{listing}[htbp]
\begin{minted}{cpp}
   // Consume the same absolute placement the simulator does: resolve every part's pose once, then
   // translate each origin-centered primitive so its fore plane lands at the resolved fore-plane
   // origin (the build* primitives are centered about mid-length, half an axialLength forward of
   // the aft plane).
   const std::vector<model::part::Placed> placed =
      model::part::resolvePlacements(*root, model::part::Pose{});
\end{minted}
\caption{The visualizer asks the model layer's resolver --- not its own arithmetic --- where
every part is. The comment names the contract. \srccap{visualizer/RocketMesh.cpp}{429}}
\label{lst:interfaces-resolver}
\end{listing}

\begin{designrule}[Rendered geometry cannot diverge from simulated geometry]
\code{resolvePlacements} is documented as ``the single authority for absolute placement:
intent (StationLink) in, geometry (Pose) out. Both the simulator and the visualizer consume
its output, so the two cannot disagree'' (\src{model/parts/Placement.h}{148}). The simulator
side consumes it inside \code{Part::ensurePlacementCache}, which feeds the composite
mass/CM/inertia pass (\src{model/parts/Part.cpp}{160}; \cref{sec:placement}); the visualizer
side is \cref{lst:interfaces-resolver}. What the user sees on screen and what the propagator
integrates are, by construction, the same poses --- there is no second placement computation
that could go stale.
\end{designrule}

Since a resolved \code{Pose::origin} is the part's \emph{fore-plane} origin and a part
occupies $z \in [-\code{axialLength}, 0]$ in its own frame (\src{model/parts/Part.h}{144}),
while the primitives are centered about mid-length, each mesh is translated by
\cppinline|p.pose.origin + Vector3(0.0, 0.0, -part.axialLength() / 2.0)|
(\src{visualizer/RocketMesh.cpp}{482}) --- the mesh center lands half an axial length aft of
the fore plane.

The overlap verdict is shared the same way. The mesh walk reads the cached
\code{SolveResult} from \code{placementDiagnostics()}
(\src{visualizer/RocketMesh.cpp}{440}); where the simulator refuses a failed solve by
throwing from the composite pass (\src{model/parts/Part.cpp}{189}), the visualizer instead
\emph{renders} the failure: every offending part is painted a fixed error red
$(0.90, 0.10, 0.10)$ regardless of the active color scheme --- the viewer ships four,
\code{Classic} (contractually first, as the default), \code{Hi-Vis}, \code{Carbon}, and
\code{Pastel} (\src{visualizer/ColorScheme.cpp}{28}) --- ``so a self-intersecting design
is unmistakable under any scheme'' (\src{visualizer/RocketGLWidget.cpp}{617}), and each
located diagnostic is also logged (\src{visualizer/RocketMesh.cpp}{446}). Same verdict
object, two consumer policies --- exactly as the \code{SolveResult} documentation prescribes
(\src{model/parts/Placement.h}{112}). The per-offender message text is carried on the render
item but not yet surfaced in the UI --- only the boolean reaches the GPU mesh
(\src{visualizer/RocketGLWidget.cpp}{515}; \cref{sec:incomplete}).

The rest of the pipeline is conventional and compact: two-sided directional lighting
(culling is off because thin fins and tube bores expose back faces,
\src{visualizer/RocketGLWidget.cpp}{305}), a stand-up transform rotating body $+z$ onto
world $+Y$ so the rocket points up on screen (\src{visualizer/RocketGLWidget.cpp}{171}), and
an orbit/pan/zoom camera that frames the model's bounding sphere on load
(\src{visualizer/RocketGLWidget.cpp}{261}). One absence completes the picture: the
visualizer has no automated test --- its CMake declares only the executable
(\srcf{visualizer/CMakeLists.txt}), no test binary references any \code{viz} symbol, and the
coverage harness aggregates neither \srcf{visualizer/} nor \srcf{gui/}
(\src{CMakeLists.txt}{248}). The poses and overlap verdicts it renders are pinned upstream
by the resolver and sweep tests (\cref{sec:testing}); the meshing, GL, and widget code paths
are eyeball-verified against the committed \code{.qrd} fixtures only.

\subsection{The maturity asymmetry}\label{sec:interfaces-asymmetry}

Read together, the three surfaces invert the usual expectation. The scriptable, headless CLI
is the flagship: complete motor management, full design authoring, a disciplined stdout
protocol, a documented flight artifact, and two test suites driving the identical object
code (\cref{sec:testing}). The desktop GUI --- the surface a newcomer would reach for first
--- is a capable motor-and-launch panel wrapped in the frame of a design tool that is not
yet there: five permanently disabled File actions, two empty panes, and no path to a
\code{.qrd}. The visualizer is complete for its narrow job and exemplary in how it stays
consistent with the simulator, but ships without a single test. The gaps named in place here
--- the disabled menu, the stub views, the missing structure-change notification, the
synchronous network and launch calls on the GUI thread --- are collected in
\cref{sec:incomplete}.
